\subsection{Goodness-of-Fit and Diagnostic Plots}

Goodness-of-fit methods compare an observed empirical pattern with the pattern predicted by a probability model.  Formal tests give calibrated evidence against a specified model, while diagnostic plots show where the model fits or fails.

\begin{definition} \label{def:multinomial-goodness-of-fit-model} \hlt{Multinomial Goodness-of-Fit Model}
\begin{enumerate}[label=(\roman*)]
\setlength{\itemsep}{0pt}
\item \hlt{Observed counts}: $O=(O_1,\ldots,O_m)$ records cell frequencies from $n$ trials, with $O_i\in\mathbb{N}_0$ and $\sum\limits_{i=1}^mO_i=n$.
\item \hlt{Null model}: $\mathcal{M}_0=\{p(\theta)=(p_1(\theta),\ldots,p_m(\theta)):\theta\in\Theta_0\}\subseteq\Delta_{m-1}$ is parametric.
\item \hlt{Saturated model}: $\Delta_{m-1}=\{p\in[0,1]^m:\sum\limits_i p_i=1\}$ allows every multinomial probability vector.
\item \hlt{Fitted counts}: if $\widehat{\theta}$ is the null MLE, then $E_i\defeq np_i(\widehat{\theta})$ for $i=1,\ldots,m$.
\item \hlt{Degrees of freedom}: if $\dim\Theta_0=k$, then $\operatorname{df}=(m-1)-k=m-k-1$.
\end{enumerate}
\end{definition}

\begin{definition} \label{def:pearson-and-likelihood-ratio-gof-statistics} \hlt{Pearson and Likelihood-Ratio Goodness-of-Fit Statistics}
\begin{enumerate}[label=(\roman*)]
\setlength{\itemsep}{0pt}
\item \hlt{Inputs}: $O=(O_1,\ldots,O_m)$ are observed counts, $E=(E_1,\ldots,E_m)$ are fitted null counts with $E_i>0$, and the null model has $k$ free parameters.
\item \hlt{Pearson statistic}: $X^2\defeq\sum\limits_{i=1}^m\frac{(O_i-E_i)^2}{E_i}$.
\item \hlt{Likelihood-ratio statistic}: $G^2\defeq2\sum\limits_{i=1}^mO_i\log(\frac{O_i}{E_i})$, with $0\log(\frac{0}{E_i})\defeq0$.
\item \hlt{Null calibration}: under the regular null model, $X^2\Rightarrow\chi^2_{m-k-1}$ and $G^2=-2\log\lambda(X)\Rightarrow\chi^2_{m-k-1}$.\\
Moreover $G^2-X^2=o_{P}(1)$ under $H_0$.
\end{enumerate}
\end{definition}

\begin{example} \label{ex:hardy-weinberg-goodness-of-fit} \hlt{Hardy--Weinberg Fit}\\
For genotype cells $(M,MN,N)$, the Hardy--Weinberg model is $p(\theta)=((1-\theta)^2,2\theta(1-\theta),\theta^2)$.\\
With $n=1029$ and $\widehat{\theta}=0.4247$, the observed and fitted counts are
\[
\begin{array}{c|ccc}
&M&MN&N\\ \hline
O_i&342&500&187\\
E_i&340.6&502.8&185.6
\end{array}
\]
Since $m=3$ and $k=1$, $\operatorname{df}=m-k-1=1$. Pearson statistic is $X^2=0.0319$, so $p=\mathbb{P}\{\chi^2_1\geq0.0319\}$ is large.\\
The data give no evidence against Hardy--Weinberg equilibrium.
\end{example}

\begin{method} \label{meth:reading-goodness-of-fit-cells} \hlt{Reading Cell Contributions}
\begin{enumerate}[label=(\roman*)]
\setlength{\itemsep}{0pt}
\item Compute signed residuals $r_i\defeq O_i-E_i$ and Pearson contributions $c_i\defeq\frac{(O_i-E_i)^2}{E_i}$.
\item Large $c_i$ identify cells driving the rejection; the sign of $r_i$ indicates overfit or underfit in that cell.
\item If some $E_i$ are small, group adjacent or tail cells before testing so the chi-square approximation is reliable.
\end{enumerate}
\end{method}

\begin{example} \label{ex:poisson-clumps-goodness-of-fit} \hlt{Poisson Clumps}\\
For grouped bacterial clump frequencies from $n=400$ grid squares, fit $H_0:X\sim\operatorname{Poisson}(\lambda)$ with $\widehat{\lambda}=2.44$ and grouped tails so $E_i$ are adequate.
The Pearson statistic is $X^2=75.4$ with $\operatorname{df}=6$, so $p=\mathbb{P}\{\chi^2_6\geq75.4\}<0.005$; the largest $c_i=\frac{(O_i-E_i)^2}{E_i}$ are in the first and last cells, indicating excess small and large counts.
\end{example}

\begin{definition} \label{def:poisson-dispersion-test} \hlt{Poisson Dispersion Test}\\
The data is $x_1,\ldots,x_n$ are independent counts and $\overline{x}\defeq n^{-1}\sum\limits_{i=1}^nx_i$.
\begin{enumerate}[label=(\roman*)]
\setlength{\itemsep}{0pt}
\item \hlt{Models}: $H_0:x_i\sim\operatorname{Poisson}(\lambda)$ with common $\lambda$; the saturated alternative uses separate rates $\lambda_i$.
\item \hlt{Likelihood-ratio statistic}: $G_D^2\defeq2\sum\limits_{i=1}^nx_i\log(\frac{x_i}{\overline{x}})$, with $0\log(\frac{0}{\overline{x}})\defeq0$.
\item \hlt{Pearson approximation}: $T\defeq\overline{x}^{-1}\sum\limits_{i=1}^n(x_i-\overline{x})^2$.
\item \hlt{Calibration}: under $H_0$, $G_D^2$ and $T$ are calibrated by $\chi^2_{n-1}$\\
Large values indicate overdispersion relative to $\operatorname{Var}(X)=\mathbb{E}X$.
\end{enumerate}
\end{definition}

\begin{example} \label{ex:bacterial-clumps-dispersion} \hlt{Dispersion in Bacterial Clumps}\\
For bacterial clump counts with $n=400$, $\overline{x}=2.44$ and $s^2=4.59$, so $\frac{s^2}{\overline{x}}=1.88>1$.\\
Dispersion statistic is $T\approx752.7$ with $\operatorname{df}=n-1=399$, so $p=\mathbb{P}\{\chi^2_{399}\geq752.7\}\approx0$.\\
The data show strong overdispersion relative to a Poisson model.
\end{example}

\begin{definition} \label{def:hanging-rootogram-and-chigram} \hlt{Hanging Rootograms and Chi-Grams}\\
Let $n_j$ be observed bin counts and $\widehat n_j$ fitted counts in ordered bins.
\begin{enumerate}[label=(\roman*)]
\setlength{\itemsep}{0pt}
\item A \hlt{hanging histogram} plots raw residuals $r_j\defeq n_j-\widehat n_j$, whose variance depends on $\widehat n_j$.
\item A \hlt{hanging rootogram} plots $q_j\defeq\sqrt{n_j}-\sqrt{\widehat n_j}$, using $\operatorname{Var}(\sqrt{n_j})\approx\frac{1}{4}$ for count data.
\item A \hlt{hanging chi-gram} plots standardized residuals $z_j\defeq\frac{n_j-\widehat n_j}{\sqrt{\widehat n_j}}$.
\end{enumerate}
Runs of same-sign residuals, tail spikes, or systematic curvature indicate model departure.
\end{definition}

\begin{method} \label{meth:probability-plot-construction} \hlt{Probability Plots}\\
Let $X_{(1)}\leq\cdots\leq X_{(n)}$ be order statistics and let $F$ be the candidate continuous strictly increasing CDF.
\begin{enumerate}[label=(\roman*)]
\setlength{\itemsep}{0pt}
\item A \hlt{P-P plot} compares $F(X_{(k)})$ with $u_k\defeq \frac{k}{n+1}$ for $k=1,\ldots,n$.
\item A \hlt{Q-Q plot} compares $X_{(k)}$ with $F^{-1}(u_k)$ for $k=1,\ldots,n$.
\end{enumerate}
The model is plausible when the plotted points are approximately linear.\\
For $F(x)=G(\frac{x-\mu}{\sigma})$, Q-Q slope and intercept estimate $\sigma$ and $\mu$.
\end{method}

\begin{method} \label{meth:probability-plots-for-grouped-data} \hlt{Probability Plots for Grouped Data}\\
Let bins boundaries $x_1<\cdots<x_{m+1}$, counts $n_1,\ldots,n_m$, total $n\defeq\sum\limits_{j=1}^mn_j$, cumulative counts $N_j\defeq\sum\limits_{i=1}^jn_i$.\\
Let $G$ be the standardised CDF in the proposed location-scale family $F_{\mu,\sigma}(x)=G(\frac{x-\mu}{\sigma})$.
\begin{enumerate}[label=(\roman*)]
\setlength{\itemsep}{0pt}
\item Plot empirical boundary quantiles $x_{j+1}$ against theoretical quantiles $G^{-1}(\frac{N_j}{n+1})$ for $j=1,\ldots,m$.
\item Use this grouped Q-Q plot when individual order statistics are unavailable.
\end{enumerate}
\end{method}

\begin{definition} \label{def:normality-diagnostics} \hlt{Normality Diagnostics}\\
Let $X_{(1)}\leq\cdots\leq X_{(n)}$, $\overline{X}\defeq n^{-1}\sum\limits_iX_i$, and $s^2\defeq n^{-1}\sum\limits_i(X_i-\overline{X})^2$.
\begin{enumerate}[label=(\roman*)]
\setlength{\itemsep}{0pt}
\item \hlt{Normal Q-Q plot}: plot $X_{(k)}$ against $\Phi^{-1}(\frac{k}{n+1})$; normality predicts approximate linearity.
\item \hlt{Skewness}: $b_1\defeq n^{-1}\sum\limits_{i=1}^n\frac{(X_i-\overline{X})^3}{s^3}$ measures asymmetry; large $\abs{b_1}$ indicates nonnormality.
\item \hlt{Kurtosis}: $b_2\defeq n^{-1}\sum\limits_{i=1}^n\frac{(X_i-\overline{X})^4}{s^4}$ measures tail weight and peakedness; normality gives $b_2\approx3$.
\item \hlt{Plot correlation}: Q-Q point correlation; small values indicate curvature.
\end{enumerate}
\end{definition}

\begin{remark} \label{rem:grouped-normality-chi-square-caution} \hlt{Grouped Normality Tests}\\
For binned normal-fit testing, use $X^2=\sum\limits_j\frac{(O_j-E_j)^2}{E_j}$ with $E_j=n\{F_{\widehat\mu,\widehat\sigma}(b_{j+1})-F_{\widehat\mu,\widehat\sigma}(b_j)\}$.\\
The usual $\chi^2_{m-k-1}$ calibration requires $(\widehat\mu,\widehat\sigma)$ from the grouped counts; ungrouped estimates change the limit.
\end{remark}
